% New English translation from checked cote 25, site 105 (PDF page 106).
% AI-translated by GPT-6 Astra (OpenAI, ChatGPT Pro mode).
\scanpage{25/105}
\mspage{8}

of $\Gamma(U,\underline{O}_U)^{\hand{*}}$.

In particular:

\uline{Corollary} 1.8. Let S be a prescheme, P=P(\underline{E}) a projective bundle over S
defined by \underline{E} loc free of finite type, u an automorphism of P. Then u
is determined, in a neighborhood of every point s $\hand{\in}$ S, by an automorphism
u of \underline{E}\hand{|}U, the latter being well determined by u modulo multiplication by
an element of $\Gamma(U,\underline{O}_U)^{\hand{*}}$.

% NOTE: the source words corresponding to "we deduce can easily deduce" are circled by hand.
\uline{Remark} 1.9. From 1.8. we \var{\xout{deduce}}{}{typewriter strike-out}{ega25_25_105-v0001} can easily deduce that the \var{\margin{\uline{do it} ! \unsure{I|i}t is the open of \unsure{$P(\underline{E}^{\vee}\otimes\underline{E})$|$P(\underline{E}\otimes\underline{E})$} defined \unsure{by|by} the nonvanishing of the ``determinant'' - - - -}}{\placedmargin{\uline{do it} ! \unsure{I|i}t is the open of \unsure{$P(\underline{E}^{\vee}\otimes\underline{E})$|$P(\underline{E}\otimes\underline{E})$} defined \unsure{by|by} the nonvanishing of the ``determinant'' - - - -}}{marginal note placed}{ega25_25_105-v0002}
\var{\xout{Aut}}{}{typewriter strike-out}{ega25_25_105-v0003} group functor $\underline{Aut}_S(P)$ on S is representable by an affine \var{\over{prescheme}{\unsure{prescheme|prescheme}}}{\unsure{prescheme|prescheme}}{writing over}{ega25_25_105-v0004}
of finite presentation over S, which is also interpreted as the \uline{quotient} group \var{\inserted[typed]{pre}}{pre}{insertion}{ega25_25_105-v0005}scheme
of the linear group scheme $\underline{Gl}(\underline{E})$ by its center
$\underline{G}_m$. This group \var{\inserted[typed]{pre}}{pre}{insertion}{ega25_25_105-v0006}scheme is called \var{\xout{classically the (prescheme}}{}{typewriter strike-out}{ega25_25_105-v0007}
the \uline{projective group prescheme}, or simply \uline{projective group}, defined
by \underline{E}, \var{\inserted[typed]{and is denoted GP(\underline{E})\hand{,}}}{and is denoted GP(\underline{E})\hand{,}}{insertion}{ega25_25_105-v0008} When \underline{E} is free, $\underline{E}\hand{\simeq}\underline{O}_S^{r+1}$, then GP(\underline{E}) is precisely the
\var{\xout{scheme}}{}{typewriter strike-out}{ega25_25_105-v0009} group prescheme $GP(r)_S$ obtained by the base change
$S\hand{\to}\mathrm{Spec}(\underline{Z})$ \var{\xout{from the}}{}{typewriter strike-out}{ega25_25_105-v0010} from the analogous group scheme GP(r) over Spec(\underline{Z}), called
the \uuline{absolute} \uline{\unsure{P|p}rojective Group}.
\tnote{25-105-1}{The uncertain overtyped letter in préschéma is r or e over the typed r; both English alternatives use prescheme, with the physical uncertainty retained. The marginal capital-or-lowercase C becomes the corresponding I or i in It, and pr or par both render by. The canceled extra é in appelé has no English spelling effect. The source uses u for both automorphisms in the corollary, and its Gl and GP notation is retained.}

\var{\margin{105}}{\placedmargin{105}}{marginal note placed}{ega25_25_105-v0012}
